\documentclass[lettersize,journal]{IEEEtran}

\usepackage{amsmath,amsfonts,amssymb}
\usepackage{algorithm}
\usepackage{algorithmic}
\usepackage{array}
\usepackage{booktabs}
\usepackage{multirow}
\usepackage[caption=false,font=footnotesize,labelfont=sf,textfont=sf]{subfig}
\usepackage{textcomp}
\usepackage{stfloats}
\usepackage{url}
\usepackage{graphicx}
\usepackage{cite}

\hyphenation{pro-ve-nance multi-scale time-line ORTHRUS}

% Compile from the project root.  All final vector figures are stored directly
% in the TDSC directory and are referenced by explicit, auditable paths.

\newcommand{\method}{CMIR}
\newcommand{\topk}{\operatorname{TopK}}
\newcommand{\rank}{\operatorname{rank}}
\newcommand{\ind}{\mathbb{I}}
\newcommand{\dataset}[1]{\textsc{#1}}

\begin{document}

\title{From Edge Anomalies to Incidents: Causal Multiscale Scoring and Fixed-Budget Reconstruction for Provenance-Based Intrusion Detection}

% Anonymous review title block.  Replace only this \author command with the
% final author/affiliation block after the journal's review-stage instructions
% permit deanonymization.  Do not add \IEEEpubid in a submission manuscript;
% IEEE supplies the publication identifier after acceptance.
\author{Anonymous Authors}

\markboth{IEEE Transactions on Dependable and Secure Computing}%
{Anonymous Authors: Causal Multiscale Scoring and Incident Reconstruction}

\maketitle

\begin{abstract}
Provenance-based intrusion detection systems score millions of system events,
whereas an analyst needs a short list of temporally coherent incidents.  This
paper presents \method{}, a post-inference layer that converts frozen anomaly
scores into ranked incidents without retraining or consulting test labels.
For each completed window, \method{} projects edge losses to role-aware node
evidence, aggregates the strongest collective evidence, and computes
forward-only scores at one- and four-window scales.  A fixed per-scale budget,
rank-percentile fusion, and deterministic temporal clustering then produce an
auditable incident queue.  Model checkpoints and reconstruction parameters are
frozen before test-label evaluation.

We evaluate five independently trained model instances on each of four
provenance benchmarks: CADETS-E3, THEIA-E3, CLEARSCOPE-E3, and OPTC-H201.  The
multiscale configuration recovers all 35 dataset--seed episode instances across
20 runs, raises mean attack-window recall from 0.607 for the efficient scale-4
mode to 0.713, and compresses 127 selected windows per cross-dataset seed into
53.8 incidents.  Its exact budget-matched joint tail probability ranges from
$1.15\!\times\!10^{-3}$ to $1.50\!\times\!10^{-3}$ across seeds.  Scale 4
produces only 29.6 incidents per cross-dataset seed while retaining 0.857
episode recall.  Ablations confirm complementary contributions from collective
evidence, temporal scale, rank fusion, and incident construction.  A node-output
adapter further composes \method{} with MAGIC, ThreaTrace, and VELOX on
CLEARSCOPE-E3; ThreaTrace and VELOX recover the attack episode in all five
seeds without detector-specific tuning.  Core reconstruction requires
5.5--5.7 ms for 161--162-window streams and 28.0 ms for a 768-window stream.
\end{abstract}

\begin{IEEEkeywords}
Advanced persistent threats, anomaly detection, forward-only temporal analysis,
incident reconstruction, multiscale fusion, provenance graphs.
\end{IEEEkeywords}

\section{Introduction}
\label{sec:introduction}

\IEEEPARstart{P}{rovenance-based} intrusion detection systems (PIDS) organize
processes, files, sockets, and system calls as causal graphs.  This representation
supports both anomaly detection and attack investigation, from early intrusion
backtracking~\cite{king2003backtracking} to streaming detectors such as
StreamSpot~\cite{manzoor2016streamspot}, UNICORN~\cite{han2020unicorn},
MAGIC~\cite{jia2024magic}, ThreaTrace~\cite{wang2022threatrace}, and
ORTHRUS~\cite{jiang2025orthrus}.  Recent systems
have substantially improved representation learning and node attribution, yet a
deployment gap remains: a detector emits edge or node scores, while a security
analyst consumes incidents.  Selecting isolated high-score events either floods
the analyst with repeated evidence from one attack or misses weak events that
become persuasive only when considered collectively and over time.

Attack reconstruction systems address related but distinct parts of this gap.
SLEUTH~\cite{hossain2017sleuth}, HOLMES~\cite{milajerdi2019holmes}, and
POIROT~\cite{milajerdi2019poirot} exploit causal dependencies to explain or align
suspicious activity; NoDoze~\cite{hassan2019nodoze} ranks alerts to reduce fatigue.
Their success motivates an explicit, testable interface between anomaly scoring
and graph investigation: which completed windows should be investigated, and
which selected windows belong to the same incident?

We answer these questions with \method{} (Causal Multiscale Incident
Reconstruction), shown in Fig.~\ref{fig:overview}.  \method{} leaves the
underlying detector unchanged.  In our experiments, ORTHRUS is trained only on
normal data, its checkpoint is selected by validation loss, and the selected
model is frozen before test inference.  \method{} then operates on timestamped
edge losses.  Its first innovation is a role-aware collective statistic: edge
losses are projected to file nodes and the strongest $K$ node values summarize
each window.  Its second innovation is strictly causal temporal evidence at two
scales.  Its third innovation is a fixed-budget rank fusion followed by temporal
incident reconstruction.  Labels enter only after the ranked incidents have
been written to disk.

ORTHRUS supplies the primary edge-output evaluation because it preserves the
full timestamped anomaly stream required by every \method{} module.  We also
compose the same frozen reconstruction parameters with node outputs from
VELOX~\cite{jiang2025orthrus}, MAGIC~\cite{jia2024magic}, and
ThreaTrace~\cite{wang2022threatrace} on CLEARSCOPE-E3.  This second experiment evaluates the
score interface directly: the upstream models are trained independently, while
\method{} receives only their node scores and native temporal locations.

Our study includes CADETS-E3, THEIA-E3, and CLEARSCOPE-E3 from DARPA Transparent
Computing~\cite{darpaTC}, together with OPTC-H201.  The four datasets differ in
operating system, activity volume, number of attack episodes, and the strength of
the frozen detector's node-level signal.  Five independent training seeds per
dataset provide 20 complete detector runs.  This design evaluates whether the
reconstruction layer remains useful across retraining variation and markedly
different score distributions.

This paper makes five contributions:
\begin{itemize}
  \item \textbf{Collective role-aware evidence.}  We transform event-level
  reconstruction loss into a robust window statistic using the strongest
  file-node evidence, preserving weak coordinated activity that a single-edge
  threshold can discard.
  \item \textbf{Strictly causal multiscale scoring.}  Short and sustained
  activity are scored with trailing windows only.  A fixed per-scale budget and
  rank-percentile fusion avoid score-scale assumptions and test-label tuning.
  \item \textbf{Incident reconstruction with an exact null model.}  Selected
  windows are merged into temporal incidents and evaluated against the exact
  uniform fixed-budget random baseline, including cross-dataset convolution.
  \item \textbf{A five-seed, four-dataset evaluation.}  The multiscale method
  recovers every annotated episode in all 20 runs, while the scale-4 operating
  mode offers a smaller analyst queue.  Complete module and sensitivity
  ablations explain the contribution of each stage.
  \item \textbf{Cross-detector composition.}  A node-output adapter evaluates
  the unchanged reconstruction parameters with three additional upstream
  detectors.  VELOX and ThreaTrace attain complete five-seed episode coverage,
  demonstrating interface compatibility beyond the primary edge-output
  implementation.
\end{itemize}

\begin{figure*}[!t]
  \centering
  \includegraphics[width=0.98\textwidth]{TDSC/fig0_method_overview.pdf}
  \caption{End-to-end overview of \method{}.  A validation-selected, frozen
  ORTHRUS model emits timestamped edge losses.  The proposed layer constructs
  collective file-node evidence, causal scores at scales 1 and 4, a fixed-budget
  rank-percentile union, and temporally coherent incidents.  Ground-truth labels
  are isolated from all stages above the dashed evaluation boundary.}
  \label{fig:overview}
\end{figure*}

\section{Background and Related Work}
\label{sec:related}

\subsection{Whole-System Provenance}

Whole-system provenance records dependencies induced by operating-system events.
Provenance-aware storage introduced persistent causal histories at the storage
layer~\cite{muniswamy2006pass}.  Linux Provenance Modules provide a trusted
whole-system capture mechanism~\cite{bates2015lpm}, while CamFlow demonstrates
practical capture, policy control, and cloud deployment
~\cite{pasquier2017camflow,pasquier2017managed}.  Kernel-supported logging
~\cite{ma2018kernel}, tamper-evident audit designs~\cite{paccagnella2020custos},
and throughput-aware auditing~\cite{inam2022faust} further strengthen the data
plane.  Provenance graphs are attractive for security because an
information-flow path naturally supports causal backtracking, but their size and
heterogeneity require strong filtering before analyst review.

\subsection{Detection and Investigation}

StreamSpot summarizes streaming heterogeneous graphs for anomaly
detection~\cite{manzoor2016streamspot}; UNICORN uses historical graph sketches
for long-running APT detection~\cite{han2020unicorn}; and MAGIC learns masked
graph representations for detection and investigation~\cite{jia2024magic}.
ThreaTrace performs node-level threat detection through provenance graph
learning~\cite{wang2022threatrace}.
SIGL learns installation-specific provenance behavior~\cite{han2021sigl},
PAGODA targets real-time detection at scale~\cite{xie2020pagoda}, and
Log2Vec learns heterogeneous enterprise-log representations
~\cite{liu2019log2vec}.  ``You Are What You Do'' uses provenance to identify
stealthy malware~\cite{wang2020yawyd}, while ATLAS learns attack sequences for
investigation~\cite{alsaheel2021atlas}.
ORTHRUS emphasizes attribution quality and emits edge reconstruction losses from
a temporal graph model~\cite{jiang2025orthrus}.  We use those frozen losses as
the upstream signal rather than modifying ORTHRUS's encoder or training loss.
This separation makes the contribution precise: \method{} is an incident-level
decision and reconstruction layer.

Investigation systems use causal structure in complementary ways.  SLEUTH
propagates trust labels over audit dependencies~\cite{hossain2017sleuth}; HOLMES
correlates suspicious information flows into APT scenarios
~\cite{milajerdi2019holmes}; POIROT aligns threat behavior with kernel audit
records~\cite{milajerdi2019poirot}; and NoDoze combines anomaly and provenance
dependence to prioritize alerts~\cite{hassan2019nodoze}.  ProTracer alternates
logging and tainting~\cite{ma2016protracer}; MPI and MCI partition or infer
causality for attack investigation~\cite{ma2017mpi,kwon2018mci}; Winnower uses
grammatical inference to compact cluster provenance~\cite{hassan2018winnower};
and RapSheet performs tactical endpoint provenance analysis
~\cite{hassan2020rapsheet}.  LogGC~\cite{lee2013loggc},
HERCULE~\cite{pei2016hercule}, and the provenance-auditing SoK
~\cite{inam2023sok} further characterize the trade-off between storage,
causality, and investigation quality.  HADES extends provenance analytics to
cross-machine Active Directory attacks in IEEE TDSC
~\cite{liu2026hades}.  \method{} supplies these downstream tools with a bounded
set of causal temporal regions, reducing repeated window alerts before expensive
graph expansion begins.

\subsection{Representation Learning and Robustness}

The frozen detector builds on a broader representation-learning lineage:
Word2Vec~\cite{mikolov2013word2vec} and node2vec~\cite{grover2016node2vec}
provide distributed and structural embeddings; GraphSAGE
~\cite{hamilton2017graphsage} supports inductive graph learning; graph attention
~\cite{velickovic2018gat} learns neighbor importance; and Temporal Graph Networks
~\cite{rossi2020tgn} model timestamped interactions.  Our rank-based fusion is
well matched to these heterogeneous score generators because it does not assume
that two temporal scales share a calibrated numeric range.  Provenance detectors
also face mimicry attacks that preserve superficially benign graph structure
~\cite{goyal2023mimicry}; collective and multiscale evidence increase the amount
of coordinated behavior that must be hidden simultaneously.

\subsection{Positioning of This Work}

Three distinctions define our approach.  First, its unit of decision is a
completed time window, not a whole-trace node maximum.  Second, all temporal
features are trailing; an earlier window score cannot change when a future event
is added.  Third, the selection budget is explicit.  The primary task is batch
triage over a closed review horizon: scores are causal at every prefix, and the
final top-$B$ selection uses the complete unlabeled score sequence in that
horizon.  These properties allow a direct operational interpretation: the
chosen budget controls the number of temporal regions admitted to incident
construction, and clustering controls the remaining incident queue.
Throughout this paper, \emph{causal} denotes forward-only temporal availability
in the systems sense; it does not denote causal-effect identification.

\section{Problem Formulation}
\label{sec:problem}

Let a provenance stream be partitioned chronologically into $T$ immutable
windows $\mathcal{W}_1,\ldots,\mathcal{W}_T$.  Window $t$ contains directed events
$e=(u,v,r,\tau)$, where $u$ and $v$ are provenance entities, $r$ is an event
type, and $\tau$ is the timestamp.  A frozen detector assigns every event a
nonnegative reconstruction loss $\ell_{e,t}$.  The reconstruction layer receives
only
\begin{equation}
  \mathcal{X}=\{(t,u,v,r,\tau,\ell_{e,t})\};
\end{equation}
it does not receive attack labels or attack intervals.  The score for window
$t$ is computed from $\mathcal{X}_{1:t}$ only; after a prespecified review
horizon closes, the unlabeled score sequence is ranked under the fixed budget.

The output is an ordered list of incidents
$\mathcal{I}=(I_1,\ldots,I_M)$, where each $I_j$ is a set of selected temporal
windows.  We seek high recall of annotated attack episodes and attack-overlap
windows under a transparent review budget.  We measure two recalls:
\begin{align}
 R_{\mathrm{win}} & = \frac{|\{t:y_t=1\}\cap\widehat{\mathcal W}|}
 {|\{t:y_t=1\}|}, \label{eq:windowrecall}\\
 R_{\mathrm{epi}} & = \frac{\sum_{q=1}^{Q}
 \ind[\widehat{\mathcal W}\cap A_q\neq\varnothing]}{Q},
 \label{eq:episoderecall}
\end{align}
where $\widehat{\mathcal W}$ is the selected-window set and $A_q$ contains the
windows overlapped by attack episode $q$.  Episode recall captures the primary
investigation objective; window recall measures within-episode temporal
coverage.

\subsection{Threat and Evaluation Model}

The adversary may distribute malicious actions across processes, files, and
time, producing evidence weaker than the largest benign event.  The adversary
does not alter stored detector losses or the experiment manifest.  Training and
checkpoint selection use no test labels.  Hyperparameters $(K,H,B,g,q)$ are
fixed at $K=10$, $H=\{1,4\}$, $B=20$ windows per scale, maximum incident gap
$g=2$, and top-score count $q=3$.  Attack annotations are consumed only in the
post-hoc evaluation represented by Eqs.~\eqref{eq:windowrecall} and
\eqref{eq:episoderecall}.

\section{Causal Multiscale Incident Reconstruction}
\label{sec:method}

\subsection{Design Objectives and End-to-End Data Flow}

The central design problem is not to produce another anomaly detector.  It is to
transform a detector's dense event-loss stream into a small number of temporal
investigation units while preserving attack evidence.  Four requirements follow
from this distinction.

\begin{enumerate}
  \item \emph{Detector preservation:} reconstruction must not change the
  trained encoder, checkpoint, edge features, or event losses.  Any improvement
  can therefore be attributed to evidence organization rather than retraining.
  \item \emph{Temporal causality:} the score attached to window $t$ may depend
  on windows $1{:}t$, but never on $t{+}1{:}T$.  This permits chronological
  replay and prevents a future event from strengthening an earlier score.
  \item \emph{Explicit workload:} admission to reconstruction must have a fixed
  and auditable budget instead of a dataset-specific score threshold.
  \item \emph{Incident semantics:} repeated evidence from one campaign should
  become one investigation item, without deleting the selected windows used to
  compute coverage.
\end{enumerate}

Figure~\ref{fig:overview} realizes these requirements as five interfaces.  The
frozen detector exports an immutable edge-loss table.  Role projection converts
edge evidence into node evidence within each window.  Collective aggregation
produces one raw score per window.  Causal temporal operators generate
scale-specific score streams, after which rank fusion enforces a per-scale
budget.  Finally, temporal clustering produces incidents and an audit manifest.
Each interface writes its output before the next stage begins; consequently,
the complete result can be reproduced from edge losses without loading the
detector checkpoint again.

The method separates two notions that are often conflated.  A \emph{causal
score} is computable when its window closes.  A \emph{budgeted selection} ranks
those causal scores over a prespecified, closed review horizon (for example, one
day or one benchmark test split).  The reported experiments evaluate this batch
triage setting.  In a deployment, the same operation can be executed at the end
of every fixed review horizon; no attack label is required in either setting.

\subsection{Why Event Thresholding Is Insufficient}

Let $N_t=|\mathcal W_t|$ be the event volume in window $t$.  A maximum edge loss
is sensitive to one benign outlier, whereas a mean edge loss can shrink as
$N_t$ grows.  A global node maximum additionally collapses time: a node that
becomes anomalous late in the trace may appear to have been suspicious in every
earlier window.  These failure modes motivate a statistic with three levels of
structure: (i) retain the strongest edge loss incident to each semantic object,
(ii) require support from several objects, and (iii) aggregate only over trailing
time.  Equations~\eqref{eq:nodeevidence}--\eqref{eq:causalscale} implement these
levels directly.

\subsection{Frozen Edge-Anomaly Signal}

For every seed, the upstream detector is trained on its prescribed normal
training split.  We select the checkpoint with minimum validation loss and
freeze both weights and neighbor state before test inference.  This protocol
prevents the reconstruction method from benefiting from a test-selected epoch.
The resulting edge losses retain fine-grained temporal and relational
information, but their raw scale varies across seeds and event types.  \method{}
therefore uses order statistics and within-seed ranks instead of a shared loss
threshold.

\subsection{Role-Aware Collective Evidence}

For a node $v$ in the chosen role $\rho$ (file in all primary experiments),
define its evidence in window $t$ as the largest incident edge loss:
\begin{equation}
 a_{v,t}=\max_{e\in\mathcal{W}_t:\,v\in\{\mathrm{src}(e),\mathrm{dst}(e)\}}
 \ell_{e,t},\qquad \operatorname{role}(v)=\rho .
 \label{eq:nodeevidence}
\end{equation}
Using both source and destination appearances is important because audit schemas
may encode the same semantic access in either direction.  Let
$a_{(1),t}\geq a_{(2),t}\geq\cdots$ be the descending role-node evidence.  The
raw collective score is
\begin{equation}
 r_t=\frac{1}{k_t}\sum_{i=1}^{k_t}a_{(i),t},\qquad
 k_t=\min(K,|V_{\rho,t}|).
 \label{eq:collective}
\end{equation}
This statistic rewards several anomalous objects without allowing window size to
dominate.  $K=1$ reduces to a maximum; averaging all nodes is
density-dominated.  The frozen primary choice $K=10$ is evaluated, rather than
selected, by the post-hoc cross-dataset ablation in Section~\ref{sec:ablation}.

\begin{figure*}[t]
  \centering
  \includegraphics[width=0.98\textwidth]{TDSC/fig10_collective_evidence_mechanism.pdf}
  \caption{Role-aware collective evidence. Edge anomaly losses are projected
  to role-constrained nodes with a max operator and then aggregated over the
  Top-$K$ node evidence. The displayed values are an explanatory example of
  Eqs.~\eqref{eq:nodeevidence}--\eqref{eq:collective}; the frozen primary
  setting is $K=10$.}
  \label{fig:collective-mechanism}
\end{figure*}

\subsection{Strictly Causal Temporal Scales}

For a trailing scale $h\in H$, \method{} computes
\begin{equation}
 c_t^{(h)}=\frac{1}{|\mathcal H_t^{(h)}|}
 \sum_{j\in\mathcal H_t^{(h)}}r_j,
 \quad
 \mathcal H_t^{(h)}=\{\max(1,t-h+1),\ldots,t\}.
 \label{eq:causalscale}
\end{equation}
Thus $h=1$ reacts to a short burst and $h=4$ accumulates sustained weak evidence.
Equation~\eqref{eq:causalscale} contains no centered or future window.  For any
$t$, $c_t^{(h)}$ is a deterministic function of $\mathcal{W}_{1:t}$, giving the
following prefix-invariance property:
\begin{equation}
 \mathcal{X}_{1:t}=\mathcal{X}'_{1:t}
 \Longrightarrow c_j^{(h)}(\mathcal{X})=c_j^{(h)}(\mathcal{X}')
 \quad \forall j\leq t.
 \label{eq:prefix}
\end{equation}

\subsection{Fixed-Budget Rank Fusion}

For each scale, rank all $T$ windows by $c_t^{(h)}$ in descending order with a
deterministic chronological tie break.  Exactly $B$ windows are selected:
\begin{equation}
 \mathcal S_h=\{t:\rank_h(t)\leq B\}.
 \label{eq:scalechoice}
\end{equation}
We map scale-specific ranks to comparable percentiles,
\begin{equation}
 p_t^{(h)}=1-\frac{\rank_h(t)-1}{\max(1,T-1)},
 \label{eq:percentile}
\end{equation}
and construct the frozen multiscale union and score
\begin{equation}
 \widehat{\mathcal W}=\bigcup_{h\in H}\mathcal S_h,
 \qquad m_t=\max_{h\in H}p_t^{(h)}.
 \label{eq:fusion}
\end{equation}
The union preserves scale-specific discoveries; the maximum percentile prevents
a strong short-lived signal from being diluted by another scale.  Because
$B|H|$ is an upper bound rather than an equality, agreement between scales
automatically reduces the review budget.

Rank normalization serves a second purpose.  If a seed, event vocabulary, or
detector checkpoint applies a strictly increasing transformation to its anomaly
losses, the ordering of windows is preserved even though their numerical ranges
differ.  This makes fusion reproducible across the markedly different score
scales observed in the four datasets.

\begin{figure*}[t]
  \centering
  \includegraphics[width=0.98\textwidth]{TDSC/fig11_causal_multiscale_fusion.pdf}
  \caption{Strictly causal multiscale scoring and rank fusion. Each temporal
  score uses only the current and preceding windows; per-scale ranks are
  normalized before the selected sets are united. Rings show an illustrative
  Top-3 selection for visual clarity, whereas all primary experiments use the
  frozen budget $B=20$ per scale and $H=\{1,4\}$.}
  \label{fig:multiscale-mechanism}
\end{figure*}

\subsection{Incident Construction and Ranking}

Sort selected window indices.  Adjacent selected indices belong to the same
incident when their difference is at most $g$; with $g=2$, one unselected window
may be bridged.  For incident $I$, let $m_{(1)}(I)\geq m_{(2)}(I)\geq\cdots$ be
its selected-window scores.  Its score is
\begin{equation}
 S(I)=\frac{1}{q_I}\sum_{i=1}^{q_I}m_{(i)}(I),
 \qquad q_I=\min(q,|I|), \quad q=3.
 \label{eq:incidentscore}
\end{equation}
Top-$q$ aggregation rewards persistent evidence but avoids penalizing a long
incident with weaker boundary windows.  Incidents are sorted by $S(I)$, then by
scale support and start time for deterministic output.

\begin{figure*}[t]
  \centering
  \includegraphics[width=0.98\textwidth]{TDSC/fig12_incident_reconstruction.pdf}
  \caption{Temporal incident reconstruction and ranking. Selected windows
  separated by at most $g=2$ are connected, and each resulting incident is
  ranked by its Top-$q$ window evidence with $q=3$. Numbers in the diagram are
  explanatory values that instantiate Eq.~\eqref{eq:incidentscore}.}
  \label{fig:incident-mechanism}
\end{figure*}

\subsection{Formal Properties}

The following properties explain why the modules compose cleanly.

\noindent\textbf{Property 1 (prefix-causal scores).}
For every $h\in H$ and $t\leq T$, $c_t^{(h)}$ is measurable with respect to the
edge losses in $\mathcal W_{1:t}$.  Appending or modifying
$\mathcal W_{t+1:T}$ cannot change $c_j^{(h)}$ for any $j\leq t$.

\emph{Justification:} $a_{v,t}$ and $r_t$ use only edges in $\mathcal W_t$.
Equation~\eqref{eq:causalscale} uses $r_j$ only for $j\leq t$.  Composition
therefore yields Eq.~\eqref{eq:prefix}.  This property is verified in the
artifact by changing a suffix and comparing every prefix score bit-for-bit.

\noindent\textbf{Property 2 (bounded review budget).}
Assume $T\geq B$ and every scale selects exactly $B$ windows.  Then
\begin{equation}
 B\leq |\widehat{\mathcal W}|\leq \min(T,B|H|).
 \label{eq:budgetbound}
\end{equation}
The lower bound occurs when all scales agree; the upper bound occurs when their
selected sets are disjoint.  Thus scale agreement automatically returns unused
capacity rather than duplicating an alert.

\noindent\textbf{Property 3 (monotone-transform invariance).}
For any scale $h$, replacing all $c_t^{(h)}$ by
$f_h(c_t^{(h)})$ for a strictly increasing function $f_h$ leaves
$\mathcal S_h$, $p_t^{(h)}$, and the multiscale union unchanged, subject to the
same deterministic tie rule.  The result follows because a strictly increasing
map preserves pairwise order.  Rank fusion therefore requires order consistency,
not cross-scale score calibration.

\noindent\textbf{Property 4 (coverage-preserving reconstruction).}
Incident clustering partitions $\widehat{\mathcal W}$ and does not discard a
selected window.  Consequently, $R_{\rm win}$ and $R_{\rm epi}$ are identical
before and after clustering, while the analyst queue satisfies
$1\leq M\leq|\widehat{\mathcal W}|$ whenever the union is nonempty.  This
separates detection coverage from presentation workload: $g$ and $q$ can change
incident boundaries or ranks without retroactively changing selected-window
recall.

\subsection{Label-Free Parameter Freezing}

The primary configuration is frozen as
\begin{equation}
 (\rho,K,H,B,g,q)=(\text{file},10,\{1,4\},20,2,3).
 \label{eq:configuration}
\end{equation}
Each parameter has an operational interpretation.  File nodes represent the
persistent objects through which many attack stages communicate; $K$ is the
amount of simultaneous object support; $H$ specifies instantaneous and sustained
evidence; $B$ is analyst capacity per scale; $g$ is temporal tolerance between
selected windows; and $q$ determines how many strong windows support an incident
rank.  The pipeline never searches these values against test labels.  The paper
reports the complete top-$K$, scale, gap, and top-$q$ grids as post-hoc ablations,
but the primary row remains Eq.~\eqref{eq:configuration} for every dataset and
seed.  Using one shared configuration is important: it tests cross-dataset
robustness rather than constructing a separate oracle for each collector.

\subsection{Efficient and Coverage-Oriented Modes}

The same evidence stream exposes two operational modes.  The \emph{efficient
mode} retains only $\mathcal S_4$, producing exactly 20 windows before incident
clustering.  It emphasizes sustained evidence and minimizes analyst workload.
The \emph{coverage mode} uses Eq.~\eqref{eq:fusion}; its union admits between 20
and 40 windows per dataset and preserves discoveries unique to either temporal
scale.  No retraining or rescaling is needed to switch modes, so an organization
can select the mode from analyst capacity rather than from hidden model
parameters.

\subsection{Node-Output Adapter}

The primary interface consumes edge losses, but the reconstruction stages only
require a timestamped set of anomaly values.  For a detector that exports one
score $z_v$ and one native temporal location $\theta_v$ per node, we form
\begin{equation}
 r_t^{\mathrm{node}}=\frac{1}{k_t}
 \sum_{i=1}^{k_t} z_{(i),t},\qquad
 \{z_{(i),t}\}=\operatorname{sort}\{z_v:\theta_v=t\},
 \label{eq:nodeadapter}
\end{equation}
where $k_t=\min(K,|\{v:\theta_v=t\}|)$ and
$z_{(1),t}\geq z_{(2),t}\geq\cdots$.  The causal scales, fixed budgets,
rank fusion, and incident construction then remain unchanged.  MAGIC and
ThreaTrace export explicit time ranges; VELOX exports the chronologically
ordered test-window index at which a node reaches its maximum loss, which is
mapped back to the immutable timestamped test-file list.  Empty native windows
receive zero evidence.  Neither node labels nor attack intervals participate in
Eq.~\eqref{eq:nodeadapter} or any subsequent ranking decision.

\begin{algorithm}[!t]
\caption{Causal Multiscale Incident Reconstruction over a Closed Review Horizon}
\label{alg:cmir}
\begin{algorithmic}[1]
\REQUIRE Ordered edge-loss windows $\mathcal{W}_{1:T}$; role $\rho$; $K$;
scales $H$; budget $B$; gap $g$; incident top-score count $q$
\FOR{$t=1$ to $T$}
  \STATE Project incident edge losses to $a_{v,t}$ using Eq.~\eqref{eq:nodeevidence}
  \STATE Compute $r_t$ from the $K$ strongest role nodes
  \FOR{each $h\in H$}
    \STATE Compute trailing score $c_t^{(h)}$ using Eq.~\eqref{eq:causalscale}
  \ENDFOR
\ENDFOR
\FOR{each $h\in H$}
  \STATE Rank $c_{1:T}^{(h)}$; store $\mathcal S_h$ and $p_t^{(h)}$
\ENDFOR
\STATE $\widehat{\mathcal W}\leftarrow\bigcup_h\mathcal S_h$
\STATE Merge ordered selected windows whose index gap is at most $g$
\STATE Score each incident using Eq.~\eqref{eq:incidentscore}
\RETURN ranked incidents and the immutable per-window audit table
\end{algorithmic}
\end{algorithm}

\subsection{Complexity and Auditability}

Let $E$ be the number of scored edges and $V_t$ the role nodes in window $t$.
Edge-to-node projection costs $O(E)$; sorting node evidence costs
$O(\sum_t |V_t|\log|V_t|)$ (or $O(\sum_t |V_t|\log K)$ with a size-$K$ heap).
Trailing means cost $O(T|H|)$ with prefix sums, scale ranking costs
$O(|H|T\log T)$, and clustering costs $O(B|H|\log(B|H|))$.  The output records
the edge-loss source, scale scores, ranks, selected scales, incident identifier,
and frozen configuration.  Every reported decision can therefore be reproduced
without rerunning detector training.

\section{Experimental Methodology}
\label{sec:setup}

\subsection{Datasets and Protocol}

Table~\ref{tab:datasets} summarizes the fixed test streams used by the incident
layer.  A window is positive when its time interval overlaps a benchmark attack
interval.  Several positive windows can belong to one attack episode; this
distinction is essential because selecting one representative window can be
sufficient to initiate investigation.

\begin{table}[!t]
\caption{Evaluation Datasets and Frozen Test Protocol}
\label{tab:datasets}
\centering
\footnotesize
\setlength{\tabcolsep}{3.6pt}
\begin{tabular}{lrrrr}
\toprule
Dataset & Test win. & Attack win. & Episodes & Seeds\\
\midrule
\dataset{CADETS-E3}     & 333 & 9 & 3 & 5\\
\dataset{CLEARSCOPE-E3} & 161 & 4 & 1 & 5\\
\dataset{THEIA-E3}      & 203 & 8 & 2 & 5\\
\dataset{OPTC-H201}     & 247 & 9 & 1 & 5\\
\bottomrule
\end{tabular}
\end{table}

All five seeds use the same data split and reconstruction hyperparameters; only
detector initialization and training randomness vary.  Validation-loss selection
chooses epochs 20, 20, 20, 19, and 19 on OPTC-H201; analogous selected
checkpoints are frozen for the other datasets.  OPTC-H201 is preprocessed into
the same 272-dimensional ORTHRUS edge representation before training.  No
ground-truth identifier or attack interval is a feature.

\subsection{Cross-Detector Composition Protocol}

We additionally train the PIDSMaker implementations of VELOX, MAGIC, and
ThreaTrace on CLEARSCOPE-E3 with five seeds each.  Every model uses its native
training and evaluation procedure and exports its epoch-11 node result.  The
adapter in Eq.~\eqref{eq:nodeadapter} uses the same
$K=10$, $H=\{1,4\}$, $B=20$, $g=2$, and $q=3$ as the primary experiment; no
detector-specific search is performed.  MAGIC and ThreaTrace expose 161 and 162
native time ranges, respectively.  VELOX exposes 768 finer-grained,
chronologically ordered edge-loss windows.  Because native temporal units
differ, episode recall and within-detector compression are the principal
composition metrics; the table also reports each native window count and
attack-window recall transparently.

\subsection{Operating Modes and Metrics}

We report two operating modes.  \emph{Scale-4} selects $B=20$ sustained-evidence
windows and is the efficient configuration.  \emph{Multiscale union} combines
scale 1 and scale 4 and is the coverage configuration.  Besides
$R_{\mathrm{win}}$ and $R_{\mathrm{epi}}$, we report selected windows, resulting
incidents, and
\begin{equation}
 C_{\mathrm{inc}}=1-\frac{M}{|\widehat{\mathcal W}|},
 \label{eq:compression}
\end{equation}
the proportional reduction from selected windows to analyst-facing incidents
after temporal merging.  Means are taken over five seeds.

The comparison set is matched to the paper's downstream scope.  With the
upstream detector and score traces fixed, we compare Top-1, Top-5, Top-10, and
all-node aggregation; one- and four-window temporal scales; their multiscale
union; and the exact uniform selector at the same realized budget.  These
comparisons test whether each reconstruction module adds value without
conflating that value with a change in detector training.  Native node-level
metrics for the external detectors are reported separately in
Table~\ref{tab:externaldetectors}.

\subsection{Exact Fixed-Budget Random Baseline}

For $T$ windows, budget $b$, and an episode occupying $a_q$ windows, the exact
probability that a uniform $b$-subset hits episode $q$ is
\begin{equation}
 P_q(\mathrm{hit})=1-\frac{\binom{T-a_q}{b}}{\binom{T}{b}}.
 \label{eq:randomhit}
\end{equation}
The expected random episode recall is
$Q^{-1}\sum_qP_q(\mathrm{hit})$.  For the tail probability of covering at least
the observed number of episodes, we count uniform fixed-size subsets by dynamic
programming over the disjoint annotated episode sets; cross-dataset
distributions are convolved under independent random selection in each dataset.
This exact null preserves each method's actual union budget and avoids a
favorable comparison to a smaller random budget.
The resulting tail probabilities quantify departure from this specified
uniform-budget selector; they are not a distribution-free guarantee for future
deployments.

\section{Results}
\label{sec:results}

\subsection{Four-Dataset Incident Recovery}

Table~\ref{tab:headline} gives the primary five-seed results.  The multiscale
union achieves 1.0 episode recall on every dataset and all 20 dataset--seed
runs.  Its mean window recall is 0.889 on CADETS, 0.950 on CLEARSCOPE, 0.800 on
THEIA, and 0.356 on OPTC.  The method therefore finds every annotated campaign
while recovering additional within-episode evidence when the frozen detector
supports it.

\begin{table*}[!t]
\caption{Five-Seed Incident Results.  ``Rnd. epi.'' Is the Expected Episode
Recall of a Uniform Random Selection With the Same Per-Seed Budget; $p$ Is the
Mean Per-Seed Exact Tail Probability.}
\label{tab:headline}
\centering
\footnotesize
\setlength{\tabcolsep}{3.8pt}
\begin{tabular}{llrrrrrrr}
\toprule
Dataset & Mode & $R_{\rm win}$ & $R_{\rm epi}$ & Sel. win. & Incidents & $C_{\rm inc}$ & Rnd. epi. & $p$\\
\midrule
\multirow{2}{*}{CADETS-E3}
 & Scale-4    & 0.778 & 0.667 & 20.0 & 7.8  & 0.610 & 0.169 & 0.071910\\
 & Multiscale & \textbf{0.889} & \textbf{1.000} & 30.4 & 14.8 & 0.513 & 0.248 & 0.012935\\
\midrule
\multirow{2}{*}{CLEARSCOPE-E3}
 & Scale-4    & \textbf{0.950} & \textbf{1.000} & 20.0 & 7.0  & \textbf{0.650} & 0.415 & 0.414885\\
 & Multiscale & \textbf{0.950} & \textbf{1.000} & 32.8 & 12.4 & 0.622 & 0.602 & 0.601833\\
\midrule
\multirow{2}{*}{THEIA-E3}
 & Scale-4    & 0.600 & \textbf{1.000} & 20.0 & 6.6  & \textbf{0.670} & 0.338 & 0.105945\\
 & Multiscale & \textbf{0.800} & \textbf{1.000} & 32.0 & 12.0 & 0.626 & 0.492 & 0.230701\\
\midrule
\multirow{2}{*}{OPTC-H201}
 & Scale-4    & 0.289 & \textbf{1.000} & 20.0 & 8.2  & \textbf{0.590} & 0.538 & 0.538421\\
 & Multiscale & \textbf{0.356} & \textbf{1.000} & 31.8 & 14.6 & 0.540 & 0.717 & 0.716645\\
\bottomrule
\end{tabular}
\end{table*}

Figure~\ref{fig:datasetrecall} visualizes this dataset-wise result.  Scale 4 is
already strong on sustained attacks: it reaches 0.95 window recall on CLEARSCOPE
and detects every episode on CLEARSCOPE, THEIA, and OPTC.  The multiscale union
adds the complementary short-scale windows needed to raise CADETS episode recall
from 0.667 to 1.0 and increases window recall by 0.20 on THEIA and 0.067 on
CADETS.  On OPTC, scale 4 raises mean attack-window recall from the scale-1 value
of 0.156 to 0.289, and fusion raises it further to 0.356.

\begin{figure*}[!t]
  \centering
  \includegraphics[width=0.88\textwidth]{TDSC/fig1_dataset_recall.pdf}
  \caption{Attack-window and attack-episode recall across four datasets.  Bars
  summarize five independently trained seeds.  Multiscale fusion attains 100\%
  episode recall on every dataset while improving or preserving temporal
  coverage relative to scale 4.}
  \label{fig:datasetrecall}
\end{figure*}

\subsection{Cross-Dataset Null Comparison and Workload}

Aggregating all four datasets within each seed, scale 4 covers six of seven
episodes ($R_{\mathrm{epi}}=0.857$), whereas multiscale covers seven of seven.
The exact budget-matched joint tail probability is $0.002306$ for scale 4.  For
multiscale it ranges from $0.001148$ to $0.001503$ because overlap between
scales changes the union budget by seed.  Figure~\ref{fig:random} compares observed and random
coverage, and Table~\ref{tab:crossdataset} reports the cross-dataset summary.

\begin{table}[!t]
\caption{Cross-Dataset Five-Seed Summary. Joint $p$ Is the Mean Exact Per-Seed
Tail Probability Under the Budget-Matched Null.}
\label{tab:crossdataset}
\centering
\footnotesize
\setlength{\tabcolsep}{3.2pt}
\begin{tabular}{lrrrrrr}
\toprule
Mode & $R_{\rm win}$ & $R_{\rm epi}$ & Sel. & Inc. & $C_{\rm inc}$ & Joint $p$\\
\midrule
Scale-4 & 0.607 & 0.857 & 80.0 & 29.6 & \textbf{0.630} & 0.002306\\
Multiscale & \textbf{0.713} & \textbf{1.000} & 127.0 & 53.8 & 0.576 & \textbf{0.001289}\\
\bottomrule
\end{tabular}
\end{table}

\begin{figure}[!t]
  \centering
  \includegraphics[width=\columnwidth]{TDSC/fig2_observed_vs_random.pdf}
  \caption{Observed episode recovery versus the exact uniform fixed-budget null.
  The advantage persists after each method is charged its actual union budget.}
  \label{fig:random}
\end{figure}

The fixed budget is also operationally meaningful.  Across four datasets,
scale 4 selects exactly 80 windows per seed and condenses them to 29.6 incidents;
multiscale selects 124--130 union windows and condenses them to 49--60 incidents.
Figure~\ref{fig:budget} shows that the additional multiscale coverage remains a
compact queue after clustering.  An operator can therefore choose the efficient
scale-4 mode or the coverage-oriented multiscale mode without changing detector
training.

\begin{figure*}[!t]
  \centering
  \includegraphics[width=0.88\textwidth]{TDSC/fig3_budget_and_compression.pdf}
  \caption{Selection budget and incident compression.  Temporal clustering
  consistently reduces the number of analyst-facing items; scale 4 gives the
  smallest queue, while multiscale spends additional budget for complete episode
  coverage.}
  \label{fig:budget}
\end{figure*}

\subsection{Stability Across Training Seeds}

The dataset-level conclusions reproduce across all five detector trainings.
For multiscale fusion, every seed recovers all seven cross-dataset episodes.
Cross-dataset attack-window recall is 21/30 or 22/30 per seed, and incident counts
range from 49 to 60.  Scale 4 covers six of seven episodes for every seed and
18--19 of 30 attack windows.  Figure~\ref{fig:stability} makes this stability
visible and confirms that the result is not driven by one favorable checkpoint.

\begin{figure}[!t]
  \centering
  \includegraphics[width=\columnwidth]{TDSC/fig4_seed_stability.pdf}
  \caption{Five-seed stability of the cross-dataset results.  Multiscale episode
  recall remains 1.0 for every independently trained model instance.}
  \label{fig:stability}
\end{figure}

\subsection{Composition With Additional Upstream Detectors}

Table~\ref{tab:externalcomposition} evaluates the node-output adapter without
changing any \method{} parameter.  ThreaTrace recovers the CLEARSCOPE attack
episode in all five seeds, reaches mean attack-window recall 0.700, and places
the attack incident at ranks 1--3.  VELOX also recovers the episode in all five
seeds.  Its finer 768-window stream yields lower within-episode window recall
and a larger queue, but demonstrates stable coverage under the same absolute
Top-20-per-scale budget.  Together, the two detectors provide ten independent
cross-detector runs with complete episode recovery.

The exact native-budget null further separates observed coverage from the
opportunity created by each detector's temporal partition.  ThreaTrace raises
episode recall from a random expectation of 0.565 to 1.0; under independent
uniform selection across the five seeded runs, the probability of covering all
five episode instances is 0.0575.  VELOX also reaches 1.0, while its broad 33-window attack
interval and 36--37-window union produce a higher random expectation of 0.806
and a five-seed tail probability of 0.339.  Thus the VELOX result establishes
interface compatibility, while ThreaTrace supplies the stronger prioritization
result.

MAGIC attains high AUC but low precision at its native operating point, and its
node-output temporal projection does not place the attack interval in the fixed
multiscale budget.  This result identifies an interface requirement: successful
fixed-budget reconstruction needs anomaly scores whose temporal projection
concentrates evidence near the attack interval; global node separability alone
is insufficient.  The unchanged configuration therefore reveals both successful
transfer and a concrete failure mode.

\begin{table*}[!t]
\caption{Five-Seed Node-Output Composition on CLEARSCOPE-E3.  ``Rnd.'' Is the
Expected Episode Recall of Uniform Selection With the Same Native Budget; $p_5$
Is the Exact Probability of Covering at Least the Observed Number of Episodes
Across Five Seeds.}
\label{tab:externalcomposition}
\centering
\footnotesize
\setlength{\tabcolsep}{3.8pt}
\begin{tabular}{lrrrrrrrrr}
\toprule
Upstream & Win. & Succ. & $R_{\rm win}$ & $R_{\rm epi}$ & Rnd. & $p_5$ & Sel. & Inc. & $C_{\rm inc}$\\
\midrule
MAGIC      & 161 & 0/5 & 0.000 & 0.000 & 0.519 & 1.0000 & 26.8 & 7.6  & 0.729\\
ThreaTrace & 162 & 5/5 & 0.700 & 1.000 & 0.565 & 0.0575 & 30.2 & 11.6 & 0.616\\
VELOX      & 768 & 5/5 & 0.036 & 1.000 & 0.806 & 0.3391 & 36.4 & 27.2 & 0.253\\
\bottomrule
\end{tabular}
\end{table*}

The best attack-incident ranks across seeds are 3, 2, 1, 2, and 2 for
ThreaTrace and 27, 20, 20, 3, and 25 for VELOX.  Absolute window and incident
counts are interpreted within each detector's native temporal resolution.

\subsection{Upstream Node-Detection Characteristics}

Table~\ref{tab:externaldetectors} reports the upstream node-level behavior
before reconstruction.  MAGIC combines high AUC with complete node recall, but
its native operating point emits 42,646 false positives on average.  ThreaTrace
also recalls every positive node and retains ranking information, while its
default threshold marks the full node population.  VELOX detects one of 41
positive nodes per seed at its native operating point.  These distinct score
profiles show that the composition experiment does not depend on a shared
node-classification operating regime; it evaluates whether the temporal score
organization remains useful to the fixed downstream interface.

\begin{table*}[!t]
\caption{Native Upstream Node-Level Performance on CLEARSCOPE-E3 (Mean $\pm$
Sample Standard Deviation Over Five Seeds).}
\label{tab:externaldetectors}
\centering
\footnotesize
\setlength{\tabcolsep}{3.7pt}
\begin{tabular}{lrrrrrrr}
\toprule
Upstream & AP & AUC & Precision & Recall & FPR & MCC & Mean FP\\
\midrule
MAGIC & $0.01654\pm0.01179$ & $0.96796\pm0.02314$ & $0.00101\pm0.00024$ & $1.000\pm0.000$ & $0.38295\pm0.10525$ & $0.02497\pm0.00503$ & 42646.4\\
ThreaTrace & $0.04181\pm0.08417$ & $0.87961\pm0.07989$ & $0.00037\pm0.00000$ & $1.000\pm0.000$ & $1.00000\pm0.00000$ & $0.00000\pm0.00000$ & 111362.0\\
VELOX & $0.000470\pm0.000033$ & $0.32953\pm0.01170$ & $0.000376\pm0.000172$ & $0.02439\pm0.00000$ & $0.04037\pm0.04578$ & $-0.00065\pm0.00284$ & 4495.2\\
\bottomrule
\end{tabular}
\end{table*}

\subsection{Reconstruction Overhead}

Table~\ref{tab:runtime} separates one-time result loading from the core
reconstruction stages.  On the 161--162-window streams, causal scoring, fusion,
and incident reconstruction finish in about 5.5--5.7 ms.  The 768-window VELOX
stream requires 28.0 ms.  Incident formation itself remains below 0.5 ms in
every configuration.  These measurements exclude upstream detector inference
and characterize the marginal reconstruction cost rather than compare detector
inference speed.  They show that post-inference reconstruction is lightweight
on the evaluated host.

\begin{table}[!t]
\caption{Post-Inference CMIR Runtime in Milliseconds (Five-Seed Mean; Core
Stages Averaged Over 100 Repetitions).}
\label{tab:runtime}
\centering
\footnotesize
\setlength{\tabcolsep}{3.0pt}
\begin{tabular}{lrrrrr}
\toprule
Upstream & Load & Score & Fusion & Incident & Core\\
\midrule
MAGIC      & 357.43 & 4.91 & 0.60 & 0.17 & $5.68\pm0.24$\\
ThreaTrace & 461.17 & 4.69 & 0.60 & 0.22 & $5.51\pm0.23$\\
VELOX      & 506.38 & 24.63 & 2.90 & 0.48 & $28.01\pm1.09$\\
\bottomrule
\end{tabular}
\end{table}

\subsection{OPTC-H201: Weak Node Signal, Strong Episode Recovery}

OPTC-H201 is the most demanding node-ranking setting.  The frozen ORTHRUS runs
contain 1,437,934 evaluated nodes and 2,905 positive nodes.  Table~\ref{tab:optcnode}
shows mean node AP of 0.0139 and recall@100 of 0.00110, yet scale 4 detects the
attack episode in every seed and multiscale also attains 1.0 episode recall.
This contrast demonstrates the value of temporal collective evidence: the method
does not require every malicious entity to occupy a high global node rank before
recovering the incident.

\begin{table}[!t]
\caption{Frozen ORTHRUS Node Ranking on OPTC-H201}
\label{tab:optcnode}
\centering
\footnotesize
\setlength{\tabcolsep}{3.8pt}
\begin{tabular}{crrrrr}
\toprule
Seed & Epoch & AP & AUC & MCC & Recall@100\\
\midrule
0 & 20 & 0.016071 & 0.738177 & 0.005754 & 0.001377\\
1 & 20 & 0.014058 & 0.736465 & 0.013940 & 0.001033\\
2 & 20 & 0.013431 & 0.736268 & 0.006920 & 0.001033\\
3 & 19 & 0.014085 & 0.737299 & 0.006078 & 0.001033\\
4 & 19 & 0.011975 & 0.735647 & 0.009211 & 0.001033\\
\midrule
Mean & -- & 0.013924 & 0.736771 & 0.008381 & 0.001102\\
SD & -- & 0.001474 & 0.000983 & 0.003388 & 0.000154\\
\bottomrule
\end{tabular}
\end{table}

\begin{figure*}[!t]
  \centering
  \subfloat[Temporal-scale gain.]{\includegraphics[width=0.47\textwidth]{TDSC/fig5_optc_scale_gain.pdf}%
  \label{fig:optcscale}}
  \hfil
  \subfloat[Frozen node-level signal.]{\includegraphics[width=0.47\textwidth]{TDSC/fig6_optc_node_signal.pdf}%
  \label{fig:optcnode}}
  \caption{OPTC-H201 analysis.  Scale 4 substantially strengthens window
  coverage relative to scale 1, while incident recovery succeeds despite a
  difficult global node-ranking problem.}
  \label{fig:optc}
\end{figure*}

\section{Ablation and Sensitivity Analysis}
\label{sec:ablation}

\subsection{Collective Evidence}

Table~\ref{tab:collective} compares the top-$K$ statistic with an all-node mean
at scale 4.  No single aggregation dominates every collector because attack
density differs across platforms.  The frozen Top-10 setting provides the most
consistent result: it recovers 91/150 attack windows and 30/35 episode instances
pooled across all dataset--seed runs.  The all-node mean recovers 28/150 windows
and 10/35 episode instances; its success on CLEARSCOPE is offset by zero episode
recovery on THEIA and OPTC.  The post-hoc ablation therefore supports Top-10 as
a robust shared setting rather than a dataset-specific oracle.

\begin{table*}[!t]
\caption{Collective-Evidence Ablation at Scale 4 (Mean $\pm$ Sample Standard
Deviation Over Five Seeds).  ``Succ.'' Counts Seeds Recovering at Least One
Attack Episode.}
\label{tab:collective}
\centering
\footnotesize
\setlength{\tabcolsep}{4.0pt}
\begin{tabular}{llrrrr}
\toprule
Dataset & Aggregation & $R_{\rm win}$ & $R_{\rm epi}$ & Succ. & Incidents\\
\midrule
\multirow{4}{*}{CADETS-E3}
 & Top-1    & $0.711\pm0.061$ & 0.667 & 5 & 7.0\\
 & Top-5    & $0.778\pm0.000$ & 0.667 & 5 & 6.8\\
 & Top-10   & $0.778\pm0.000$ & 0.667 & 5 & 7.8\\
 & All-mean & $0.178\pm0.061$ & 0.333 & 5 & 5.6\\
\midrule
\multirow{4}{*}{CLEARSCOPE-E3}
 & Top-1    & $0.000\pm0.000$ & 0.000 & 0 & 5.8\\
 & Top-5    & $0.000\pm0.000$ & 0.000 & 0 & 7.0\\
 & Top-10   & $0.950\pm0.112$ & 1.000 & 5 & 7.0\\
 & All-mean & $1.000\pm0.000$ & 1.000 & 5 & 5.6\\
\midrule
\multirow{4}{*}{THEIA-E3}
 & Top-1    & $0.800\pm0.068$ & 1.000 & 5 & 4.6\\
 & Top-5    & $0.625\pm0.000$ & 1.000 & 5 & 6.4\\
 & Top-10   & $0.600\pm0.056$ & 1.000 & 5 & 6.6\\
 & All-mean & $0.000\pm0.000$ & 0.000 & 0 & 5.4\\
\midrule
\multirow{4}{*}{OPTC-H201}
 & Top-1    & $0.156\pm0.099$ & 0.800 & 4 & 9.0\\
 & Top-5    & $0.178\pm0.061$ & 1.000 & 5 & 8.8\\
 & Top-10   & $0.289\pm0.061$ & 1.000 & 5 & 8.2\\
 & All-mean & $0.000\pm0.000$ & 0.000 & 0 & 8.4\\
\bottomrule
\end{tabular}
\end{table*}

\begin{figure}[!t]
  \centering
  \includegraphics[width=\columnwidth]{TDSC/fig7_collective_evidence_ablation.pdf}
  \caption{Collective-evidence ablation.  Sparse top-$K$ aggregation is markedly
  more robust than averaging the full active-node population.}
  \label{fig:collective}
\end{figure}

\subsection{Multiscale Fusion}

Figure~\ref{fig:fusionablation} compares individual scales and their union.
Scale 4 increases CLEARSCOPE window recall from 0.45 to 0.95 relative to scale 1,
raises CADETS window recall from 0.711 to 0.778, and nearly doubles OPTC window
recall from 0.156 to 0.289.  Scale 1 remains complementary: its union with scale
4 restores the third CADETS attack episode and raises THEIA window recall from
0.60 to 0.80.  This explains why the efficient and coverage modes occupy
different useful operating points rather than representing redundant parameter
settings.

\begin{figure*}[!t]
  \centering
  \includegraphics[width=0.88\textwidth]{TDSC/fig8_multiscale_fusion_ablation.pdf}
  \caption{Temporal-fusion ablation across scales 1 and 4.  Sustained evidence
  supplies the largest single-scale gain, while the union recovers complementary
  short-lived behavior and all attack episodes.}
  \label{fig:fusionablation}
\end{figure*}

\subsection{Incident Reconstruction Sensitivity}

The final module has two interpretable controls: maximum gap $g$ and top-score
count $q$.  Figure~\ref{fig:sensitivity} reports the complete grid over four
datasets and five seeds.  Increasing $g$ merges nearby evidence and reduces the
review queue; $q$ controls whether an incident rank reflects its peak or several
supporting windows.  The frozen choice $(g,q)=(2,3)$ consistently compresses the
selected union while retaining the episode coverage established before
clustering.  Because clustering never removes a selected window, incident
construction preserves $R_{\mathrm{win}}$ and $R_{\mathrm{epi}}$ by design.

\begin{figure*}[!t]
  \centering
  \includegraphics[width=0.88\textwidth]{TDSC/fig9_incident_reconstruction_sensitivity.pdf}
  \caption{Incident-reconstruction sensitivity to temporal gap and incident
  score aggregation.  The primary setting $g=2$, $q=3$ offers a stable balance
  between incident compression and concise ranking.}
  \label{fig:sensitivity}
\end{figure*}

\section{Discussion}
\label{sec:discussion}

\subsection{Operational Interpretation}

\method{} provides two deployable choices from the same frozen detector.
Scale 4 is appropriate when analyst capacity is tight: it admits 20 windows per
dataset and compresses them by 63\% in the cross-dataset aggregate.  Multiscale
fusion is appropriate when campaign coverage is primary: it achieves complete
episode recall across all 20 runs and still compresses the union by 57.6\%.
Ranks, selected scales, and incident membership are explicit, allowing an
analyst to understand why an item entered the queue.
The node-output experiment adds a complementary deployment path: the same
reconstruction parameters recover the CLEARSCOPE episode in every VELOX and
ThreaTrace seed, allowing \method{} to operate behind an existing scorer rather
than replacing it.  Under the exact native-budget null, ThreaTrace improves
episode recall from an expectation of 0.565 to 1.0, with five-seed tail
probability 0.0575.  VELOX provides consistent interface-level episode coverage
under a much finer temporal partition, while its larger attack opportunity is
made explicit by the corresponding null expectation of 0.806.

\subsection{Why the Gains Are Complementary}

The experiments identify three distinct effects.  Top-$K$ aggregation suppresses
the mass of ordinary active objects while preserving coordinated evidence.
Causal scale 4 accumulates weak behavior that is not individually extreme.
Scale 1 adds sharp transient behavior, and incident clustering converts both
signals into analyst-facing units.  The combined method succeeds on OPTC even
when global node ranking is weak, while retaining near-complete temporal coverage
on CLEARSCOPE.  This diversity is the central advantage of the reconstruction
layer.

\subsection{Scope and Future Deployment}

The primary edge-output study uses four public provenance benchmarks with five
independent trainings per dataset, and the composition study adds fifteen
CLEARSCOPE runs from three upstream detector architectures.  These experiments
provide evidence across operating systems, attack structures, retraining
variation, and detector interfaces.  The composition experiment establishes
compatibility with three score formats on CLEARSCOPE; extending this controlled
comparison across datasets is a natural next step.  Likewise,
OPTC contains one annotated attack episode spread across nine windows; its
episode metric is therefore interpreted together with window recall and the
exact budget-matched null.  A production deployment can apply the same score
interface at the end of fixed review horizons, but live analyst studies and
non-DARPA environments remain future work.  Runtime measurements depend on the
host and storage stack;
we therefore separate one-time loading from the repeatedly measured core
stages.  Across the external streams, core reconstruction requires only
5.51--28.01 ms, and incident formation itself remains below 0.5 ms.

\section{Conclusion}
\label{sec:conclusion}

This paper introduced \method{}, a strictly causal multiscale scorer with a
fixed-budget reconstruction layer that turns
frozen edge anomalies into a compact incident queue.  Role-aware collective
evidence, trailing temporal scales, rank-percentile fusion, and deterministic
incident clustering address complementary failure modes without changing the
underlying detector.  Across CADETS-E3, CLEARSCOPE-E3, THEIA-E3, and OPTC-H201,
the multiscale method recovered all 35 dataset--seed episode instances in 20
independently trained runs and improved mean attack-window recall to 0.713.  The
mean exact cross-dataset budget-matched tail probability was approximately
$1.3\times10^{-3}$,
and clustering reduced 127 selected windows to 53.8 incidents per seed.  The
efficient scale-4 mode provides a second strong operating point with a 29.6-item
queue and 0.857 episode recall.  Without detector-specific tuning, the
node-output adapter additionally recovers the CLEARSCOPE episode in all five
VELOX and all five ThreaTrace runs.  ThreaTrace improves on the exact random
episode-recall expectation of 0.565 and attains a five-seed tail probability of
0.0575; the budget-matched VELOX null transparently quantifies its broader
attack opportunity.  Core reconstruction completes in 5.51--28.01 ms across
the three external streams.  On the evaluated traces, these results show that
forward-only collective scoring and fixed-budget reconstruction convert
high-volume provenance anomaly streams into reproducible,
investigation-ready incident queues.

\section*{Artifact Availability}

The artifact records the selected checkpoint for every seed, immutable edge-loss
directories, scale-specific window tables, multiscale union manifests, incident
membership, post-hoc label evaluations, exact random-budget calculations, and
the external-detector adapter manifests and outputs.  External artifacts include
per-seed detector metrics, component runtimes over 100 repetitions, and the
native-budget random-baseline tables.  The artifact also includes the scripts
that regenerate all thirteen figures and paper tables.  The regression suite
includes a prefix-invariance test: changing a future window cannot change any
earlier raw or causal score.

\begin{thebibliography}{99}

\bibitem{king2003backtracking}
S.~T. King and P.~M. Chen, ``Backtracking intrusions,'' in \emph{Proc. 19th ACM
Symp. Operating Systems Principles (SOSP)}, 2003, pp. 223--236.

\bibitem{manzoor2016streamspot}
E.~Manzoor, S.~M. Milajerdi, and L.~Akoglu, ``Fast memory-efficient anomaly
detection in streaming heterogeneous graphs,'' in \emph{Proc. 22nd ACM SIGKDD
Int. Conf. Knowledge Discovery and Data Mining (KDD)}, 2016, pp. 1035--1044.

\bibitem{han2020unicorn}
X.~Han, T.~Pasquier, A.~Bates, J.~Mickens, and M.~Seltzer, ``UNICORN: Runtime
provenance-based detector for advanced persistent threats,'' in \emph{Proc. Netw.
Distrib. Syst. Security Symp. (NDSS)}, 2020.

\bibitem{jia2024magic}
Z.~Jia, Y.~Xiong, Y.~Nan, Y.~Zhang, J.~Zhao, and M.~Wen, ``MAGIC: Detecting
advanced persistent threats via masked graph representation learning,'' in
\emph{Proc. 33rd USENIX Security Symp.}, 2024, pp. 5197--5214.

\bibitem{wang2022threatrace}
S.~Wang, Z.~Wang, T.~Zhou, H.~Sun, X.~Yin, D.~Han, H.~Zhang, X.~Shi, and
J.~Yang, ``ThreaTrace: Detecting and tracing host-based threats in node level
through provenance graph learning,'' \emph{IEEE Trans. Inf. Forensics Security},
vol.~17, pp. 3972--3987, 2022, doi: 10.1109/TIFS.2022.3208815.

\bibitem{jiang2025orthrus}
B.~Jiang, T.~Bilot, N.~El Madhoun, K.~Al Agha, A.~Zouaoui, S.~Iqbal, X.~Han,
and T.~Pasquier, ``ORTHRUS: Achieving high quality of attribution in
provenance-based intrusion detection systems,'' in \emph{Proc. 34th USENIX
Security Symp.}, 2025.

\bibitem{hossain2017sleuth}
M.~N. Hossain, S.~M. Milajerdi, J.~Wang, B.~Eshete, R.~Gjomemo, R.~Sekar,
S.~Stoller, and V.~N. Venkatakrishnan, ``SLEUTH: Real-time attack scenario
reconstruction from COTS audit data,'' in \emph{Proc. 26th USENIX Security
Symp.}, 2017, pp. 487--504.

\bibitem{milajerdi2019holmes}
S.~M. Milajerdi, B.~Eshete, R.~Gjomemo, and V.~N. Venkatakrishnan, ``HOLMES:
Real-time APT detection through correlation of suspicious information flows,''
in \emph{Proc. IEEE Symp. Security and Privacy (S\&P)}, 2019, pp. 1137--1152.

\bibitem{milajerdi2019poirot}
S.~M. Milajerdi, R.~Gjomemo, B.~Eshete, R.~Sekar, and V.~N. Venkatakrishnan,
``POIROT: Aligning attack behavior with kernel audit records for cyber threat
hunting,'' in \emph{Proc. ACM SIGSAC Conf. Computer and Communications Security
(CCS)}, 2019, pp. 1795--1812.

\bibitem{hassan2019nodoze}
W.~U. Hassan, S.~Guo, D.~Li, Z.~Chen, K.~Jee, Z.~Li, and A.~Bates, ``NoDoze:
Combatting threat alert fatigue with automated provenance triage,'' in
\emph{Proc. Netw. Distrib. Syst. Security Symp. (NDSS)}, 2019.

\bibitem{bates2015lpm}
A.~Bates, D.~J. Tian, K.~R.~B. Butler, and T.~Moyer, ``Trustworthy whole-system
provenance for the Linux OS,'' in \emph{Proc. 24th USENIX Security Symp.}, 2015,
pp. 319--334.

\bibitem{pasquier2017camflow}
T.~Pasquier, J.~Han, T.~Moyer, A.~Bates, O.~Hermant, D.~Eyers, J.~Bacon, and
M.~Seltzer, ``Practical whole-system provenance capture,'' in \emph{Proc. ACM
Symp. Cloud Computing (SoCC)}, 2017, pp. 405--418.

\bibitem{muniswamy2006pass}
K.-K. Muniswamy-Reddy, D.~A. Holland, U.~Braun, and M.~Seltzer,
``Provenance-aware storage systems,'' in \emph{Proc. USENIX Annu. Tech. Conf.
(ATC)}, 2006, pp. 43--56.

\bibitem{pasquier2017managed}
T.~F.~J.-M. Pasquier, J.~Singh, D.~Eyers, and J.~Bacon, ``CamFlow: Managed
data-sharing for cloud services,'' \emph{IEEE Trans. Cloud Computing}, vol.~5,
no.~3, pp. 472--484, 2017.

\bibitem{ma2018kernel}
S.~Ma, J.~Zhai, Y.~Kwon, K.~H. Lee, X.~Zhang, G.~Ciocarlie, A.~Gehani,
V.~Yegneswaran, D.~Xu, and S.~Jha, ``Kernel-supported cost-effective audit
logging for causality tracking,'' in \emph{Proc. USENIX Annu. Tech. Conf.
(ATC)}, 2018, pp. 241--254.

\bibitem{paccagnella2020custos}
R.~Paccagnella, P.~Datta, W.~U. Hassan, A.~Bates, C.~W. Fletcher, A.~Miller,
and D.~Tian, ``Custos: Practical tamper-evident auditing of operating systems
using trusted execution,'' in \emph{Proc. Netw. Distrib. Syst. Security Symp.
(NDSS)}, 2020.

\bibitem{inam2022faust}
M.~A. Inam, A.~Goyal, J.~Liu, J.~Mink, N.~Michael, S.~Gaur, A.~Bates, and
W.~U. Hassan, ``FAuST: Striking a bargain between forensic auditing's security
and throughput,'' in \emph{Proc. 38th Annu. Computer Security Applications Conf.
(ACSAC)}, 2022, pp. 813--826.

\bibitem{han2021sigl}
X.~Han, X.~Yu, T.~Pasquier, D.~Li, J.~Rhee, J.~Mickens, M.~Seltzer, and
H.~Chen, ``SIGL: Securing software installations through deep graph learning,''
in \emph{Proc. 30th USENIX Security Symp.}, 2021, pp. 2345--2362.

\bibitem{xie2020pagoda}
Y.~Xie, D.~Feng, Y.~Hu, Y.~Li, S.~Sample, and D.~D.~E. Long, ``PAGODA: A
hybrid approach to enable efficient real-time provenance based intrusion
detection in big data environments,'' \emph{IEEE Trans. Dependable Secure
Comput.}, vol.~17, no.~6, pp. 1283--1296, 2020.

\bibitem{liu2019log2vec}
F.~Liu, Y.~Wen, D.~Zhang, X.~Jiang, X.~Xing, and D.~Meng, ``Log2Vec: A
heterogeneous graph embedding based approach for detecting cyber threats within
enterprise,'' in \emph{Proc. ACM SIGSAC Conf. Computer and Communications
Security (CCS)}, 2019, pp. 1777--1794.

\bibitem{wang2020yawyd}
Q.~Wang, W.~U. Hassan, D.~Li, K.~Jee, X.~Yu, K.~Zou, J.~Rhee, Z.~Chen,
W.~Cheng, C.~A. Gunter, and H.~Chen, ``You are what you do: Hunting stealthy
malware via data provenance analysis,'' in \emph{Proc. Netw. Distrib. Syst.
Security Symp. (NDSS)}, 2020.

\bibitem{alsaheel2021atlas}
A.~Alsaheel, Y.~Nan, S.~Ma, L.~Yu, G.~Walkup, Z.~B. Celik, X.~Zhang, and
D.~Xu, ``ATLAS: A sequence-based learning approach for attack investigation,''
in \emph{Proc. 30th USENIX Security Symp.}, 2021, pp. 3005--3022.

\bibitem{ma2016protracer}
S.~Ma, X.~Zhang, and D.~Xu, ``ProTracer: Towards practical provenance tracing
by alternating between logging and tainting,'' in \emph{Proc. Netw. Distrib.
Syst. Security Symp. (NDSS)}, 2016.

\bibitem{ma2017mpi}
S.~Ma, J.~Zhai, F.~Wang, K.~H. Lee, X.~Zhang, and D.~Xu, ``MPI: Multiple
perspective attack investigation with semantic aware execution partitioning,''
in \emph{Proc. 26th USENIX Security Symp.}, 2017, pp. 1111--1128.

\bibitem{kwon2018mci}
Y.~Kwon, F.~Wang, W.~Wang, K.~H. Lee, W.-C. Lee, S.~Ma, X.~Zhang, D.~Xu,
S.~Jha, G.~Ciocarlie, A.~Gehani, and V.~Yegneswaran, ``MCI: Modeling-based
causality inference in audit logging for attack investigation,'' in \emph{Proc.
Netw. Distrib. Syst. Security Symp. (NDSS)}, 2018.

\bibitem{hassan2018winnower}
W.~U. Hassan, N.~Aguse, M.~Lemay, T.~Moyer, and A.~Bates, ``Towards scalable
cluster auditing through grammatical inference over provenance graphs,'' in
\emph{Proc. Netw. Distrib. Syst. Security Symp. (NDSS)}, 2018.

\bibitem{hassan2020rapsheet}
W.~U. Hassan, A.~Bates, and D.~Marino, ``Tactical provenance analysis for
endpoint detection and response systems,'' in \emph{Proc. IEEE Symp. Security
and Privacy (S\&P)}, 2020, pp. 1172--1189.

\bibitem{lee2013loggc}
K.~H. Lee, X.~Zhang, and D.~Xu, ``LogGC: Garbage collecting audit log,'' in
\emph{Proc. ACM SIGSAC Conf. Computer and Communications Security (CCS)}, 2013,
pp. 1005--1016.

\bibitem{pei2016hercule}
K.~Pei, Z.~Gu, B.~Saltaformaggio, S.~Ma, F.~Wang, Z.~Zhang, L.~Si, X.~Zhang,
and D.~Xu, ``HERCULE: Attack story reconstruction via community discovery on
correlated log graph,'' in \emph{Proc. 32nd Annu. Computer Security Applications
Conf. (ACSAC)}, 2016, pp. 583--595.

\bibitem{inam2023sok}
M.~Inam, Y.~Chen, A.~Goyal, J.~Liu, J.~Mink, N.~Michael, S.~Gaur, A.~Bates,
and W.~U. Hassan, ``SoK: History is a vast early warning system: Auditing the
provenance of system intrusions,'' in \emph{Proc. IEEE Symp. Security and Privacy
(S\&P)}, 2023, pp. 307--325.

\bibitem{goyal2023mimicry}
A.~Goyal, X.~Han, A.~Bates, and G.~Wang, ``Sometimes, you aren't what you do:
Mimicry attacks against provenance graph host intrusion detection systems,'' in
\emph{Proc. Netw. Distrib. Syst. Security Symp. (NDSS)}, 2023.

\bibitem{mikolov2013word2vec}
T.~Mikolov, I.~Sutskever, K.~Chen, G.~S. Corrado, and J.~Dean, ``Distributed
representations of words and phrases and their compositionality,'' in \emph{Proc.
Advances in Neural Information Processing Systems (NeurIPS)}, 2013,
pp. 3111--3119.

\bibitem{grover2016node2vec}
A.~Grover and J.~Leskovec, ``node2vec: Scalable feature learning for networks,''
in \emph{Proc. 22nd ACM SIGKDD Int. Conf. Knowledge Discovery and Data Mining
(KDD)}, 2016, pp. 855--864.

\bibitem{hamilton2017graphsage}
W.~L. Hamilton, R.~Ying, and J.~Leskovec, ``Inductive representation learning
on large graphs,'' in \emph{Proc. Advances in Neural Information Processing
Systems (NeurIPS)}, 2017, pp. 1024--1034.

\bibitem{velickovic2018gat}
P.~Veli\v{c}kovi\'c, G.~Cucurull, A.~Casanova, A.~Romero, P.~Li\`o, and
Y.~Bengio, ``Graph attention networks,'' in \emph{Proc. Int. Conf. Learning
Representations (ICLR)}, 2018.

\bibitem{rossi2020tgn}
E.~Rossi, B.~Chamberlain, F.~Frasca, D.~Eynard, F.~Monti, and M.~Bronstein,
``Temporal graph networks for deep learning on dynamic graphs,''
arXiv:2006.10637, 2020.

\bibitem{liu2026hades}
Q.~Liu, K.~Bao, W.~U. Hassan, and V.~Hagenmeyer, ``HADES: Detecting and
investigating Active Directory attacks via whole network provenance analytics,''
\emph{IEEE Trans. Dependable Secure Comput.}, vol.~23, no.~1, pp. 936--953,
2026.

\bibitem{darpaTC}
Defense Advanced Research Projects Agency, ``Transparent Computing,'' program
and evaluation datasets, 2015--2019. [Online]. Available:
\url{https://www.darpa.mil/research/programs/transparent-computing}

\end{thebibliography}

\end{document}
